\subsection{Integrantes del equipo}
El desarrollo de SHIFT se realizó mediante un equipo de trabajo organizado en diferentes roles y responsabilidades, cada integrante participó en una parte específica del desarrollo de la aplicación web, permitiendo distribuir las actividades de análisis, diseño y programación.
\begin{itemize}
\item \textbf{Luis David Balandrano Delgado -- Líder técnico:} encargado de coordinar el desarrollo técnico del proyecto, organizar las actividades del equipo y supervisar la integración de los diferentes componentes de la aplicación
\item \textbf{Keydhy Mareli Villa Aldana -- Analista:} encargada del análisis y documentación del proyecto, así como de la elaboración y organización de los requerimientos, documentación técnica y evidencias del desarrollo
\item \textbf{Isis Dayane Arellano Carillo -- Frontend:} encargada del desarrollo de interfaces relacionadas con el diagnóstico inicial y el dashboard principal del usuario
\item \textbf{Vaneza Lugo Ramírez -- Backend:} encargada del desarrollo de funcionalidades relacionadas con las tarjetas de vocabulario y la biblioteca o sección de repaso
\item \textbf{Brian Jael Juárez Rodríguez -- Frontend:} encargado del desarrollo de las interfaces correspondientes a los ejercicios interactivos de aprendizaje
\item \textbf{Julio  Cesar Licona Acosta -- Backend:} encargado del desarrollo de funcionalidades relacionadas con el mapa de lecciones y las recomendaciones de la siguiente lección
\end{itemize}

\subsection{Distribución de interfaces}
La distribución de las interfaces se realizó de acuerdo con las actividades asignadas a cada integrante mediante el sistema de Issues del repositorio de GitHub
\begin{itemize}
\item \textbf{Keydhy -- Selección de áreas de aprendizaje:} desarrolló la pantalla donde el usuario selecciona el área de inglés que desea aprender. La interfaz incluye las opciones de Programación, Gastronomía, Turismo, Negocios e Inglés cotidiano, además de la selección visual y la validación del botón de continuar.
\item \textbf{Isis -- Diagnóstico inicial:} desarrolló la pantalla del diagnóstico inicial, destinada a evaluar el nivel de inglés del usuario mediante preguntas, selección de respuestas y presentación del resultado final.
\item \textbf{Isis -- Dashboard y progreso:} desarrolló el dashboard principal del alumno, incluyendo elementos visuales para mostrar el progreso general, la racha de días y gráficas de avance.
\item \textbf{Brian -- Ejercicios interactivos:} desarrolló las vistas de ejercicios para practicar el idioma mediante actividades como completar frases, arrastrar palabras, selección múltiple y emparejamiento.
\item \textbf{Vaneza -- Tarjetas de vocabulario:} desarrolló las tarjetas de aprendizaje o flashcards, incluyendo la palabra en inglés, traducción, definición, pronunciación y ejemplos.
\item \textbf{Vaneza -- Biblioteca y repaso:} desarrolló la pantalla de repaso y glosario, con una organización visual del vocabulario aprendido y opciones de búsqueda y filtros.
\item \textbf{Julio -- Mapa de lecciones:} desarrolló la vista de árbol o ruta de aprendizaje, donde se representan las lecciones desbloqueadas, bloqueadas y completadas.
\item \textbf{Julio 
-- Recomendación de siguiente lección:} desarrolló la sección destinada a mostrar una recomendación de la siguiente lección dentro del dashboard.
\item \textbf{Luis David -- Autenticación:} desarrolló las interfaces relacionadas con el registro de usuario, inicio de sesión, recuperación de contraseña, verificación por correo y perfil.
\item \textbf{Luis David -- Resultados y calificación:} desarrolló la pantalla de finalización de la lección, donde se presenta la calificación, puntuación y resumen de aciertos y errores.
\end{itemize}
